\section*{Author Sections 6.4--6.5: positivity of middle matrices}
Write $Q_c^R=Q_{c,0}^R-Q_{c,1}^R+\cdots\pm Q_{c,d}^R$ in Lemma6.19's approximate decomposition. The following is the key lemma, proved through the next subsections.
\paragraph{Lemma 6.21.} W.p. $>1-n^{-8\log n}$: (1) for all $R\in\binom{[n]}{\le d/2}$,
\[Q_{0,0}^R-Q_{0,1}^R+\cdots\pm Q_{0,d/2}^R\succeq\tau^{-7\tau}\operatorname{diag}(\widetilde{\mathrm{Cl}})_{\mathcal S^R\times\mathcal S^R};\]
(2) for all $R$, $0<c\le|R|$,
\[\pm\omega^{-c}(Q_{c,0}^R-Q_{c,1}^R+\cdots\pm Q_{c,d/2}^R)\preceq n^{-c/4}\operatorname{diag}(\widetilde{\mathrm{Cl}})_{\mathcal S^R\times\mathcal S^R}.\]
The proof has three ingredients: Corollary6.36 shows positivity of $Q_{0,0}^R$ via a structural/pseudorandom decomposition; Lemmas6.37/6.38 bound the rest. The last real technical challenge is $\mathbb E[Q_{0,0}^R]$, ultimately the choice of $F$ in Definition3.11.
\paragraph{Definition 6.22.} The root diagonal-clique matrix on $\binom{[n]}{\le d/2}$ is
\[D_{\mathrm{Cl}}(A,B)=\begin{cases}0,&A\ne B,\\2^{-\binom{|A|}2/2}\widetilde{\mathrm{Cl}}_A=2^{-\binom{|A|}2/2}\sum_{T\subseteq E[A]}\chi_T,&A=B.\end{cases}\tag{6.40}\]
Thus $D_{\mathrm{Cl}}^2(A,A)=\widetilde{\mathrm{Cl}}(A)$. Define $D_{\mathrm{Cl}}^\tau:=D_{\mathrm{Cl}}\otimes\mathrm{Id}_{\{0,\ldots,\tau\}\times\{0,\ldots,\tau\}}$ (6.41), also diagonal.
\paragraph{Definition 6.23.} The structural-pseudorandom decomposition is
\[Q_{0,0}^R=D_{\mathrm{Cl}}^\tau\mathbb E[Q_{0,0}^R]D_{\mathrm{Cl}}^\tau+\left(Q_{0,0}^R-D_{\mathrm{Cl}}^\tau\mathbb E[Q_{0,0}^R]D_{\mathrm{Cl}}^\tau\right).\tag{6.42}\]
The first summand is structural and the second pseudorandom.
\paragraph{Proposition 6.24 (Expression of $\mathbb E[Q_{c,0}^R]$).} Fix $R$, $0\le c\le|R|$, $r=|R|$.
\begin{enumerate}\item The expectation is supported on blockwise partial diagonals $\{((A,i),(A,j))\in\mathcal S^R\times\mathcal S^R\}$, requiring $R\subseteq A$ and $|A|+\min\{i,j\}\ge d/2$.
\item On these positions,
\begin{align*}\mathbb E[Q_{c,0}^R]((A,i),(A,j))&=\sum_{l=c}^r(-1)^{r-l}\frac{\binom rl}{(l-c)!}\binom{|A|+i+j+l-d}c\\&\quad\cdot\frac{(|A|+8\tau^2+(l-c)+(i+j))!}{(8\tau^2)!}+O(\tau^{1.5\tau}/n).\tag{6.43}\end{align*}
For $c=0$ this is
\[\sum_{l=0}^r(-1)^{r-l}\frac{\binom rl}{l!}\frac{(|A|+8\tau^2+l+(i+j))!}{(8\tau^2)!}+O(\tau^{1.5\tau}/n).\tag{6.44}\]
\item If $1_{A,A}$ is the matrix with one 1 at $(A,A)$, then
\[\mathbb E[Q_{0,0}^R]=\sum_{\substack{A\in\binom{[n]}{\le d/2}\\A\supseteq R}}1_{A,A}\otimes\left[\sum_{l=0}^r(-1)^{r-l}\frac{\binom rl}{l!}P_{|A|+l}+E_A^R\right].\tag{6.45}\]
$P_{|A|+l}$ and $E_A^R$ are $(\tau+1)\times(\tau+1)$ matrices supported on $\{i:d/2-|A|\le i\le\tau\}^2$,
\[\|E_A^R\|<\tau^{2\tau}/n,\tag{6.46}\]
and
\[P_{|A|+l}(i,j)=\frac{(|A|+l+8\tau^2+(i+j))!}{(8\tau^2)!}.\tag{6.47}\]
\end{enumerate}
\begin{proof}(1) Constant terms in(6.23) have $T_m=\emptyset$, nonzero only when $A=B$. (2) For those positions the expectation is $Y_c(r,a)$, $a=|A|+i+j$, namely
\[\sum_{l=c}^r(-1)^{r-l}\binom rl\binom{a+l-d}c\binom{n-a}{l-c}n^{-(l-c)}\frac{(a+l-c+8\tau^2)!}{(8\tau^2)!}.\tag{6.48}\]
Use
\[\binom{n-a}{l-c}n^{-(l-c)}=\frac1{(l-c)!}\frac{(n-a)\cdots(n-a-(l-c)+1)}{n^{l-c}}=\frac1{(l-c)!}(1-O(d^2/n)),\]
and the bound
\[\binom rl\binom{a+l-d}c\binom{n-a}{l-c}n^{-(l-c)}\frac{(a+l-c+8\tau^2)!}{(8\tau^2)!}<(4d)^d(9\tau^2)^d<\tau^\tau\]
to obtain(6.43). For $c=0$, the author writes that $\binom{a+l-d}0=0$ regardless of $a+l-d$ in(6.48), and the same analysis gives(6.44). (3) Each error matrix has dimension $(\tau+1)^2$ and entries absolutely $<\tau^{1.5\tau}/n$ by(2).\end{proof}
\paragraph{Remark 6.25.} The specialty of $c=0$ in(6.44) is that $\binom{|A|+l-d}0$ is always 1. For $c>0$, $\binom{|A|+l-d}c$ can be zero or negative depending on the order between $0,c,|A|+l-d$, so its expected matrix may not be PSD.
\paragraph{Definition 6.26.} For $m,t\in\mathbb N$, the factorial Hankel matrix is $H_{m,t}(i,j)=(i+j+t)!$ for $0\le i,j\le m$ (6.49).
\paragraph{Proposition 6.27 (Almost common decomposition).} (1)
\[H_{m,t}=L_m[N_{m,t}D_{m,t}(N_{m,t})^\top]L_m^\top,\]
where $L_m,D_{m,t}$ are diagonal and $N_{m,t}$ lower-triangular,
\[L_m(i,i)=i!,\quad D_{m,t}(i,i)=\prod_{t'=1}^t(i+t'),\quad N_{m,t}(i,j)=\binom{i+t}{i-j}.\]
$L_m$ is independent of $t$ and $H_{m,t}$ positive. (2) Let $J_m$ be the $(m+1)\times(m+1)$ lower-triangular Jordan block, equal to 1 if $i=j$ or $i=j+1$ and 0 otherwise. Then $N_{m,t+1}=N_{m,t}J_m$ (6.50).
\begin{proof}This follows from a direct inspection.\end{proof}
\paragraph{Proposition 6.28.} If $t+1>8\max\{r^2,m\}$ (6.51), then $H_{m,t+1}\succeq2r^2H_{m,t}$.
\begin{proof}By Proposition6.27 it suffices to show $J_mD_{m,t+1}J_m^\top\succeq2r^2D_{m,t}$. For $x^\top=(x_0,\ldots,x_m)$ compare
\[q_{t+1}(x):=(x^\top J_m)D_{m,t+1}(J_m^\top x),\qquad q_t(x):=2r^2x^\top D_{m,t}x.\tag{6.52}\]
Put $\alpha(y)=2r^2\prod_{t'=1}^t(y+t')$, $\beta(y)=\prod_{t'=1}^{t+1}(y+t')$. Then $q_{t+1}(x)=\sum_{i=0}^m\beta(i)(x_i+x_{i+1})^2$, $x_{m+1}:=0$, while $q_t(x)=\sum_{i=0}^m\alpha(i)x_i^2$. Completing squares gives
\begin{multline*}q_{t+1}(x)=\sum_{i=0}^m\left[\alpha(i)x_i^2+(\beta(i)-\alpha(i))\left(x_i+\frac{\beta(i)}{\beta(i)-\alpha(i)}x_{i+1}\right)^2\right.\\\left.-\frac{\beta(i)^2}{\beta(i)-\alpha(i)}x_{i+1}^2\right].\end{multline*}
Define $b_0=1-\alpha(0)/\beta(0)$ and
\[b_i=1-\frac{\alpha(i)}{\beta(i)}-\frac{\beta(i-1)}{\beta(i)}\frac1{b_{i-1}}.\tag{6.53}\]
Then
\[q_{t+1}(x)=\underbrace{\sum_{i=0}^m\alpha(i)x_i^2}_{q_t(x)}+\sum_{i=0}^m\beta(i)b_i\left(x_i+\frac1{b_i}x_{i+1}\right)^2.\tag{6.54}\]
\paragraph{Claim 6.29.} Each $b_i>1/2$ for $i\le m$.
\paragraph{Proof of claim.} $b_0=1-2r^2/(t+1)$ and, for $i\ge1$,
\[b_i=1-\frac{2r^2}{t+1+i}-\frac i{t+1+i}\frac1{b_{i-1}}.\tag{6.55}\]
Induct: $b_0>1/2$ by(6.51), and
\[b_i\ge1-\frac{2r^2}{t+1}-\frac m{t+1}\cdot2>1/2\]
by(6.51) and the inductive hypothesis. This proves the claim. Since $b_i$ are positive, (6.54) gives $q_{t+1}\ge q_t$, proving the proposition.\end{proof}
Apply this to(6.47): on its support,
\[P_{|A|+l}=\frac1{(8\tau^2)!}H_{\tau-(d/2-|A|),\ d-|A|+8\tau^2+l}.\]
\paragraph{Corollary 6.30 (Positiveness of $\mathbb E[Q_{0,0}^R]$).} In(6.45), regarding matrices on their support,
\[\sum_{l=0}^r(-1)^{r-l}\frac{\binom rl}{l!}P_{|A|+l}+E_A^R\succ\operatorname{diag}(\tau^{-6\tau})_{0\le i\le\tau-(d/2-|A|)}.\tag{6.56}\]
In particular,
\[\mathbb E[Q_{0,0}^R]\succ\sum_{A\supseteq R}1_{A,A}\otimes\operatorname{diag}(\tau^{-6\tau})_{d/2-|A|\le i\le\tau}=\operatorname{diag}(\tau^{-6\tau})_{\mathcal S^R\times\mathcal S^R}.\tag{6.57}\]
\begin{proof}The last part follows by checking support and that tensoring with a nonzero PSD matrix preserves $\succ$. For(6.56), fix $A$ and put
\[\tau_0=\tau-(d/2-|A|),\qquad t_0=d-|A|+8\tau^2.\tag{6.58}\]
Then
\[\sum_{l=0}^r(-1)^{r-l}\frac{\binom rl}{l!}P_{|A|+l}=\frac1{(8\tau^2)!}(X_r+X_{r-2}+\cdots),\tag{6.59}\]
where for $0\le v\le\lfloor r/2\rfloor$,
\[X_{r-2v}=\frac{\binom r{r-2v}}{(r-2v)!}\left(H_{\tau_0,t_0+r-2v}-\frac{(r-2v)^2}{2v+1}H_{\tau_0,t_0+r-2v-1}\right),\quad H_{\tau_0,-1}:=0.\]
The second coefficient is $\le r^2$. As $t_0>8\max\{r^2,\tau_0\}$, Proposition6.28 gives
\[X_{r-2v}\succeq\frac{\binom r{r-2v}}{(r-2v)!}\max\{\tfrac12H_{\tau_0,t_0+r-2v},\ r^2H_{\tau_0,t_0+r-2v-1}\}.\]
In particular,
\begin{align*}\sum_{l=0}^r(-1)^{r-l}\frac{\binom rl}{l!}P_{|A|+l}&\succeq\frac1{(8\tau^2)!}H_{\tau_0,t_0}\\&=L[N_{t_0}\tfrac{D_{t_0}}{(8\tau^2)!}(N_{t_0})^\top]L^\top,\tag{6.60}\end{align*}
omitting index $\tau_0$ on the RHS. Claim6.31 below implies the RHS is $\succ L\operatorname{diag}(\tau^{-5\tau})L\succeq\operatorname{diag}(\tau^{-5\tau})$. By Proposition6.24(3), $\|E_A^R\|<\tau^{2\tau}/n<\tau^{-6\tau}$ in the parameter regime. Thus the LHS of(6.56) is $\succeq\operatorname{diag}(\tau^{-5\tau}-\tau^{-6\tau})\succeq$ its RHS.\end{proof}
\paragraph{Claim 6.31.} In Corollary6.30's notation,
\[N_{t_0}^{-1}(i,j)=(-1)^{i-j}\binom{i+t_0}{i-j},\quad0\le i,j\le\tau_0,\tag{6.61}\]
and
\[N_{t_0}\frac{D_{t_0}}{(8\tau^2)!}(N_{t_0})^\top\succ\operatorname{diag}(\tau^{-5\tau})_{0\le i\le\tau_0}.\tag{6.62}\]
\begin{proof}For(6.61), multiplying by $N_{t_0}$ gives at $(i,j)$
\begin{align*}\sum_{j\le k\le i}(-1)^{i-k}\binom{i+t_0}{i-k}\binom{k+t_0}{k-j}&=\sum_{k'=0}^{i'}(-1)^{i'-k'}\binom{i'+j+t_0}{i'-k'}\binom{k'+j+t_0}{k'},\end{align*}
where $i'=i-j$, $k'=k-j$. To see identity, let $D_m[(1+x)^a]$ be the coefficient of $x^m$ for $m\ge0,a\in\mathbb Z$. The RHS is
\begin{align*}&(-1)^{i'}\sum_{k'=0}^{i'}D_{i'-k'}[(1+x)^{i'+j+t_0}]D_{k'}[(1+x)^{-(t_0+j+1)}]\\&=(-1)^{i'}D_{i'}[(1+x)^{i'+j+t_0-(t_0+j+1)}]\\&=(-1)^{i'}D_{i'}[(1+x)^{i'-1}]=\mathbf1_{i'=0}.
\end{align*}
For(6.62), equivalently
\[\frac{D_{t_0}}{(8\tau^2)!}\succ N_{t_0}^{-1}\tau^{-5\tau}(N_{t_0}^{-1})^\top.\tag{6.63}\]
Set $a_0=\tau^{-5\tau}$. The RHS quadratic form is
\[x^\top N_{t_0}^{-1}a_0(N_{t_0}^{-1})^\top x=a_0\sum_{j=0}^{\tau_0}y_j^2,\tag{6.64}\]
with $y_j=(x^\top N_{t_0}^{-1})_j=\sum_{i=j}^{\tau_0}(-1)^{i-j}\binom{i+t_0}{i-j}x_i$. By Cauchy--Schwarz, $y_j^2\le\tau_0\sum_{i=j}^{\tau_0}\binom{i+t_0}{i-j}^2x_i^2$. Hence
\begin{align*}a_0\sum_jy_j^2&\le a_0\sum_{i=0}^{\tau_0}x_i^2\left(\tau_0\sum_{j=0}^i\binom{i+t_0}{i-j}^2\right)\\&<\sum_{i=0}^{\tau_0}\tau^{-5\tau}(9\tau^2)^{2i+2}x_i^2.\end{align*}
But the LHS diagonal of(6.63) satisfies
\[\frac{D_{t_0}(i,i)}{(8\tau^2)!}\ge(8\tau^2)^{-(d/2-|A|)}>\tau^{-2d}>\tau^{-5\tau}(9\tau^2)^{2i+2},\]
using $i\le\tau_0<\tau$ and $d\ll\tau$. Thus(6.63) holds.\end{proof}
\paragraph{Corollary 6.32 (Positive structural part).}
\[D_{\mathrm{Cl}}^\tau\mathbb E[Q_{0,0}^R]D_{\mathrm{Cl}}^\tau\succeq\tau^{-6\tau}\operatorname{diag}(\widetilde{\mathrm{Cl}})_{\mathcal S^R\times\mathcal S^R}.\]
\begin{proof}This follows from Corollary6.30 and $D_{\mathrm{Cl}}^2(A,A)=\widetilde{\mathrm{Cl}}(A)$ in Definition6.22.\end{proof}
